%% kapa: exam=1 round=32 year=2021 period=1 form=A subject=civil_law subject_ko=민법 seq=29 number=29
%% source: Q-Net 감정평가사 기출문제 (cst00309) artlSeq=5208397 file=extracted/1008/5208397/A형/2021년도 제32회 감정평가사 1차 1교시 A형.pdf
%% conversion: kordoc 4.24.51 -> tools/gen1.py (auto); review=proofread-sonnet; flags=none
\begin{question}{29}
공유관계에 관한 설명으로 옳지 않은 것은? (다툼이 있으면 판례에 따름)
\begin{choices}
  \item 부동산 공유자의 공유지분 포기의 의사표시가 다른 공유자에게 도달하더라도 이로써 곧바로 공유지분 포기에 따른 물권변동의 효력이 발생하는 것은 아니다.
  \item 소수지분권자는 공유물의 전부를 협의 없이 점유하는 다른 소수지분권자에게 공유물의 인도를 청구할 수 있다.
  \item 과반수 지분권자는 공유물의 관리에 관한 사항을 단독으로 결정할 수 있다.
  \item 토지공유자 사이에서는 지분비율로 공유물의 관리비용을 부담한다.
  \item 공유자는 특별한 사정이 없는 한 언제든지 공유물의 분할을 청구할 수 있다.
\end{choices}
\end{question}
\answer{2}
